%% ================= TRANG BÌA =================
%% ================= TRANG BÌA =================
\hypersetup{pageanchor=false}
\begin{titlepage}
\newgeometry{margin=0cm}
\begin{tikzpicture}[remember picture, overlay]
\begin{scope}[shift={(current page.south west)}]

  % ---------- LỚP 1: NỀN LAVENDER + LƯỚI Ô VUÔNG ----------
  \fill[lavender] (0,0) rectangle (21,29.7);
  \begin{scope}[line width=0.3pt, cyanaccent!35!white]
    \foreach \x in {0,0.9,...,21} { \draw (\x,0) -- (\x,29.7); }
    \foreach \y in {0,0.9,...,29.7} { \draw (0,\y) -- (21,\y); }
  \end{scope}
  \begin{scope}[line width=0.6pt, cyanaccent!55!white]
    \foreach \x in {0,3.6,...,21} { \draw (\x,0) -- (\x,29.7); }
    \foreach \y in {0,3.6,...,29.7} { \draw (0,\y) -- (21,\y); }
  \end{scope}

  % ---------- LỚP 2: ĐƯỜNG CONG LỚN XUYÊN TRANG ----------
  \begin{scope}[opacity=0.15]
    \fill[cyan] plot[domain=1:20, samples=60, smooth]
      (\x,{6+13*(1-((\x-10.5)*(\x-10.5))/90.25)}) -- (20,4) -- (1,4) -- cycle;
  \end{scope}
  \draw[navy, line width=1.3pt, opacity=0.18, domain=1:20, samples=60, smooth]
    plot (\x,{6+13*(1-((\x-10.5)*(\x-10.5))/90.25)});

  % ---------- LỚP 3: KHỐI TIÊU ĐỀ 3D ----------
  % Nới rộng khung: (3.4, 11.7) đến (17.6, 21.2)
  \fill[navy!85!black, opacity=0.35, rounded corners=3pt]
    ($(3.4,11.7)+(0.35,-0.35)$) rectangle ($(17.6,21.2)+(0.35,-0.35)$);
  \draw[cyan, line width=1.6pt, rounded corners=3pt]
    ($(3.4,11.7)+(0.15,-0.15)$) rectangle ($(17.6,21.2)+(0.15,-0.15)$);
  \filldraw[fill=white, draw=navy, line width=2.3pt, rounded corners=3pt]
    (3.4,11.7) rectangle (17.6,21.2);

  \node[text=navy, font=\sansbold\fontsize{32}{36}\selectfont]
    at (10.5,19.1) {ỨNG DỤNG HÌNH HỌC};
  \node[text=navy, font=\sansbold\fontsize{32}{36}\selectfont]
    at (10.5,17.2) {CỦA TÍCH PHÂN};
  \draw[cyan, line width=1pt] (6.7,15.8) -- (14.3,15.8);
  \node[text=navy!80!black, font=\sffamily\slshape\fontsize{13}{16}\selectfont]
    at (10.5,14.5) {Chuyên đề: Diện tích hình phẳng \quad$\bullet$\quad Thể tích vật thể};

  % ---------- LỚP 4: HAI HÌNH MINH HỌA LỚN (Đã phóng to & dời vào giữa) ----------
  % Hình trái: Diện tích hình phẳng (Đã thêm trục Oxy, cận a, b)
  \node[opacity=0.85] at (5.5, 8.5) {%
    \begin{tikzpicture}[scale=0.55]
      \fill[cyan, opacity=0.4]
         (2,2.5) .. controls (4,6) and (6,6) .. (8,3.5)
         .. controls (6,1.5) and (4,0.5) .. (2,2.5) -- cycle;
      
      \draw[-Latex, thick, navy] (-0.5,0) -- (10.5,0) node[below right] {$x$};
      \draw[-Latex, thick, navy] (0,-0.5) -- (0,6.5) node[above left] {$y$};
      \node[below left, navy] at (0,0) {$O$};

      \draw[navy, thick]
         (0.5,0) .. controls (0.8,0.5) and (1.2,1.1) .. (2,2.5)
                 .. controls (4,6) and (6,6) .. (8,3.5)
                 .. controls (8.8,2.5) and (9.2,2.0) .. (9.5,1.5);
      \draw[teal, thick]
         (0.5,4.5) .. controls (0.8,4.0) and (1.2,3.3) .. (2,2.5)
                   .. controls (4,0.5) and (6,1.5) .. (8,3.5)
                   .. controls (8.8,4.3) and (9.2,4.7) .. (9.5,5.0);
                   
      \draw[dashed, navy, opacity=0.8, thick] (2,0) node[below, font=\bfseries] {$a$} -- (2,2.5);
      \draw[dashed, navy, opacity=0.8, thick] (8,0) node[below, font=\bfseries] {$b$} -- (8,3.5);
    \end{tikzpicture}%
  };

  % Hình phải: Khối tròn xoay (Đã phóng to, thêm trục Oxy)
  \node[opacity=0.85] at (15.5, 8.5) {%
    \begin{tikzpicture}[declare function={f(\x)=sqrt(\x);}]
      \begin{axis}[
          hide axis,
          xmin=-0.5, xmax=5.5, ymin=-2.8, ymax=3.2,
          width=6.5cm, height=4.8cm,
          clip=false, samples=60]
          
        \draw[-Latex, thick, navy] (-0.5,0) -- (5.2,0) node[below right] {$x$};
        \draw[-Latex, thick, navy] (0,-2.5) -- (0,3.0) node[above left] {$y$};
        \node[below left, navy] at (0,0) {$O$};

        \fill[cyan, fill opacity=0.22]
          plot[domain=0:4, smooth] (\x,{f(\x)}) -- (4,-2)
          -- plot[domain=4:0, smooth] (\x,{-f(\x)}) -- cycle;
        \fill[cyan, fill opacity=0.32] (4,0) ellipse (0.35 and 2);
        
        \draw[navy, thick] (4,-2) arc (-90:90:0.35 and 2);
        \draw[navy, thick, dashed] (4,2) arc (90:270:0.35 and 2);
        \draw[navy, thick, opacity=0.4] (2.5,-1.58) arc (-90:90:0.22 and 1.58);
        \draw[navy, thick, dashed, opacity=0.4] (2.5,1.58) arc (90:270:0.22 and 1.58);
        
        \addplot[domain=0:4, navy, thick, smooth] {f(x)};
        \addplot[domain=0:4, navy, thick, smooth] {-f(x)};
      \end{axis}
    \end{tikzpicture}%
  };

  % Khung công thức (Nằm ngay dưới 2 hình)
  \node[draw=navy, draw opacity=0.6, rounded corners=3pt, fill=white, fill opacity=0.9,
        text=navy, inner sep=5pt, font=\normalsize]
    at (5.5, 5.0) {$S=\displaystyle\int_a^b\ababs{f(x)-g(x)}\dx$};

  \node[draw=navy, draw opacity=0.6, rounded corners=3pt, fill=white, fill opacity=0.9,
        text=navy, inner sep=5pt, font=\normalsize]
    at (15.5, 5.0) {$V=\pi\displaystyle\int_a^b f^2(x)\dx$};

  % ---------- LỚP 5: THÔNG TIN TÁC GIẢ (Đã hạ xuống dưới) ----------
  \node[text=navy, font=\sffamily\bfseries\fontsize{14}{16}\selectfont]
    at (10.5, 3.5) {Biên soạn: Thầy Tâm dạy Toán};
  \node[text=navy!80!black, font=\sffamily\fontsize{12}{14}\selectfont]
    at (10.5, 2.7) {Liên hệ (Zalo): 0837715745};

  % ---------- LỚP 6: THANH TRÊN -- DÒNG BẮT BUỘC ----------
  \fill[navy] (0,28.2) rectangle (21,29.7);
  \node[text=white, font=\sffamily\bfseries\fontsize{11}{13}\selectfont]
    at (10.5,28.95) {CHƯƠNG TRÌNH GIÁO DỤC PHỔ THÔNG 2018 -- MÔN TOÁN 12};

  % ---------- LỚP 7: HUY HIỆU "TOÁN" + QUẢ CẦU VÀNG "12" ----------
  \foreach \i in {1,...,8} { % chktex 11
    \node[anchor=west, text=navybg, opacity=0.5,
          font=\condbold\bfseries\fontsize{74}{74}\selectfont]
      at ($(4.6,24.3)+(\i*0.014,-\i*0.014)$) {TOÁN};
  }
  \node[anchor=west, text=navytext, font=\condbold\bfseries\fontsize{74}{74}\selectfont]
    at (4.6,24.3) {TOÁN};

  % Đã dời quả cầu sang trái (x=14.0) và hạ thấp xuống (y=24.4) để cân bằng
  \fill[black, opacity=0.10] ($(15.5,24)+(0.25,-0.25)$) circle (1.8);
  \shade[ball color=gold] (15.5,24) circle (1.65);
  \fill[white, opacity=0.28] ($(15.5,24)+(-0.53,0.53)$) circle (0.65);
  \node[text=navybg, font=\sansbold\fontsize{44}{44}\selectfont] at (15.5,24) {12};
  % ---------- LỚP 8: THANH DƯỚI -- PHIÊN BẢN ----------
  \fill[navy] (0,0) rectangle (21,1.5);
  \node[text=white, font=\sffamily\fontsize{10}{12}\selectfont]
    at (10.5,0.75) {Phiên bản 1.0.0 \ $\vert$ \ Tháng 9/2026};

  \end{scope}
\end{tikzpicture}
\end{titlepage}
\hypersetup{pageanchor=true}
\restoregeometry
